\documentclass[12pt]{article}
\usepackage[T1]{fontenc}
\usepackage[utf8]{inputenc}
\usepackage{lmodern}
\usepackage{textcomp}
\usepackage{enumitem}
\usepackage{geometry}
\usepackage{graphicx}
\usepackage{amsmath, amssymb}
\geometry{margin=1in}
\begin{document}
\begin{center}
Psychology - Undergraduate Level Assessment 17 \
Total Marks: 40 \
Answer all questions.
\end{center}

\section*{Multiple Choice (10 marks)}
\begin{enumerate}[label=\arabic*.]
\item Verbal stimuli may be an acquired source of drive. Discuss.
\begin{enumerate}[label=(\alph*)]
    \item Verbal cues are effective as motivational sources only in specific cultural contexts.
    \item Only non-verbal stimuli can serve as sources of drive.
    \item Verbal stimuli serve as a source of drive exclusively in early childhood development.
    \item Verbal stimuli are an innate source of drive, not acquired.
    \item Verbal stimuli are not an acquired source of drive.
\end{enumerate}
\item In Pavlov's conditioning of dogs, the point at which the dogs salivated at the sound of the tone without the food being present is referred to as
\begin{enumerate}[label=(\alph*)]
    \item reinforcement
    \item an unconditioned stimulus
    \item conditioned response
    \item operant conditioning
    \item generalization
\end{enumerate}
\item The belief that a child's misbehavior has one of four goals - i.e., attention, revenge, power, or to display inadequacy - is most consistent with:
\begin{enumerate}[label=(\alph*)]
    \item Skinner's behaviorism
    \item Adler's individual psychology.
    \item Beck's cognitive-behavioral therapy.
    \item Perls's Gestalt therapy.
    \item Erikson's stages of psychosocial development
\end{enumerate}
\item Discuss the effects of sleep deprivation.
\begin{enumerate}[label=(\alph*)]
    \item Decreased blood pressure, heightened sensory perception, faster reaction times
    \item Strengthened cognitive abilities, longer lifespan, better pain tolerance
    \item Enhanced problem-solving skills, increased emotional stability, improved coordination
    \item Sharper focus, accelerated learning ability, increased creativity
    \item Improved cardiovascular health, reduced inflammation, increased energy levels
\end{enumerate}
\item What function does time have in classical andoperant conditioning ?
\begin{enumerate}[label=(\alph*)]
    \item Intensity of the unconditioned stimulus
    \item Conditioned stimuli
    \item Temporal contiguity
    \item Sequence of stimuli
    \item Spatial contiguity
\end{enumerate}
\end{enumerate}

\section*{True / False (10 marks)}
\textit{Answer True or False}
\begin{enumerate}[label=\arabic*.]
\setcounter{enumi}{5}
\item True or False: A group's performance on an additive task is limited by the skills of the least knowledgeable member.
\item True or False: The focus of structuralists aligns most closely with the behaviorists' emphasis on observable behavior.
\item True or False: Greater belief change occurs when compliance is accompanied by a large reward in forced compliance situations.
\item True or False: The correction for attenuation formula is used to assess the impact of increasing the number of test takers on test validity.
\item True or False: Men are less likely than women to experience stress in crowded conditions.
\end{enumerate}

\section*{Long Form Response (20 marks)}
\textit{Answer each question with a clear and complete explanation.}
\begin{enumerate}[label=\arabic*.]
\setcounter{enumi}{10}
\item Discuss how Rosenhan's research with "pseudopatients" illustrates the impact of environmental factors on the perception of behavior within psychiatric settings.
\item Discuss the ethical implications of a therapist's personal feelings toward a client, using Dr. Leonard Lykowski's situation as a case study. What steps should he take to address this issue?
\end{enumerate}

\vfill
\noindent\textit{End of Paper}
\end{document}